\section{Introduction}

\IEEEPARstart{U}{rban} traffic congestion imposes severe and compounding economic, environmental, and public-health costs on rapidly growing metropolises. Adaptive traffic signal control (TSC) seeks to relieve this burden by allocating right-of-way in response to real-time demand rather than to a fixed schedule. Over the past decade, deep reinforcement learning (RL) has become the dominant paradigm for adaptive TSC, advancing from single-intersection Deep Q-Networks~\cite{wei2018_intellilight,genders2016_drl}
to pressure-based and cooperative multi-agent formulations~\cite{wei2019_presslight,wei2019_colight,chen2020_mplight}, and most recently to centralized-training, decentralized-execution (CTDE) multi-agent actor--critic methods that coordinate large urban
networks~\cite{yu2022_mappo,mao2023_masac,zhou2024_micdrl,gu2024_partition}. Across this lineage, learned controllers consistently outperform classical fixed-time and actuated heuristics \emph{in simulation}~\cite{wei2019_survey,rltsc_review2025,rltsc_survey2026}. Yet a critical deployment bottleneck persists: almost none of these controllers can be operated on the commodity roadside sensing infrastructure that actually exists at signalized intersections.

\subsection{The Observation Gap, Not the Dynamics Gap}
The principal barrier to deployment is what we term the \emph{observation gap}.
Prevailing RL-TSC frameworks consume simulator-privileged state---exact halting vehicle counts, per-vehicle accumulated waiting times, and network-wide arrival profiles---quantities that a microscopic simulator exposes for free but that no roadside camera can recover at inference time. We deliberately distinguish this
\emph{observation gap} from the \emph{dynamics gap} (discrepancies in the underlying physics or transition delays) that dominates sim-to-real transfer in robotics~\cite{peng2018_dynrand,zhao2020_sim2real_survey} and that recent TSC transfer studies also target~\cite{da2024_promptgat}: the two are orthogonal, and closing one does not close the other. Legacy inductive-loop detectors offer only a partial remedy, since point sensing captures presence but not spatial composition, vehicle class, or continuous speed profiles~\cite{zhang2021_partialdetection}. Camera-based RL-TSC has begun to narrow the state side of the gap by extracting occupancy or queue features from
video~\cite{dai2022_worldmodels,azfar2025_cosim,su2025_visionllm}; however, these works constrain only the \emph{observation} space while still training and evaluating against a \emph{privileged reward} (e.g., simulator delay or network throughput) that is unavailable at the edge. The paired restriction of
\emph{both} state and reward to camera-computable signals therefore remains substantially under-addressed. This paper closes the observation gap by construction and quantifies its restriction cost in simulation; closed-loop physical field operation is reported in a companion study and is explicitly outside the present scope.

\subsection{Challenges}
Engineering an RL-TSC controller that factors entirely through commodity sensing raises three coupled technical challenges.
\emph{(i) Spatial granularity.} A camera region of interest (ROI) measures traffic as an approach-level aggregate and cannot supply the per-lane occupancy indices that classical state formulations presuppose, particularly where lane discipline is weak.
\emph{(ii) Reward unmeasurability.} The most effective RL reward signals---system-wide delay and exact per-vehicle waiting time---are not computable at the edge, so a reward optimized in simulation generally cannot be reproduced at deployment without a privileged read.
\emph{(iii) Heterogeneous traffic dynamics.} In developing metropolises, traffic is heavily motorcycle-dominant and non-lane-based: motorcycles filter and swarm between lanes~\cite{ngoc2021_fivecities,vn_roundabout}, a behavior that 
breaks car-centric, lane-indexed state and reward designs and that demands continuous-space (sublane) microsimulation rather than discrete lane queues.

\subsection{Contributions}
We propose a deployability-constrained multi-agent proximal policy optimization (MAPPO) framework that addresses these challenges through four verifiable contributions.
\begin{itemize}
  \item \textbf{C1 (Camera-computable state proxy).} We design a
  $26$-dimensional per-agent observation built on per-\emph{approach}
  aggregation rather than per-lane indexing. Each approach contributes bounded features (effective queue, occupancy, average speed, motorcycle share, heavy-vehicle share) and the controller state is completed by a phase one-hot vector, a normalized green timer, and a \emph{signed group-mean} pressure scalar. Every feature admits a documented camera estimator from a YOLOv11~\cite{khanam2024_yolov11}
  plus ByteTrack~\cite{zhang2022_bytetrack} pipeline, specified in
  Section~\ref{sec:obs_space} (Table~\ref{tab:obs_space}); its residual recovery error is
  bounded by the sensing-noise model of Section~\ref{sec:noise_model}, with full per-feature
  calibration against labeled ground truth deferred to the companion field study.
  \item \textbf{C2 (Proxy-computable reward).} We formulate a
  \emph{queue-centric} vision-proxy reward
  in which \emph{every} term factors through the same camera projection as the state---no term reads a privileged simulator quantity, the accumulated-delay penalty included. Its congestion core pairs the $L_2$ norm of the per-approach halted-queue signal (a mean-of-squares penalty that discourages starvation via Jensen's inequality) with the $L_1$ norm of the \emph{same} signal (a linearly-averaged delay proxy); this core is regularized by a signed \emph{phase-imbalance} pressure term---an intra-junction demand gradient, distinct from network max-pressure---that supplies a reward gradient even when allocation is correct, and by a local ROI-exit throughput term whose simulator computation is semantically identical to ByteTrack track terminations at an ROI boundary.
  \item \textbf{C3 (Information-restriction cost).} We isolate and quantify the \emph{information-restriction cost} by comparing the proxy-constrained policy against an otherwise-identical privileged-state policy evaluated on pristine features, and we characterize graceful degradation under a sensing-noise model whose magnitudes are set on a documented basis from the detection-and-tracking literature---with the class-confusion rate measured from the detector's own confusion matrix---rather than being arbitrarily assumed (Section~\ref{sec:noise_model}).
  \item \textbf{C4 (Mixed-traffic, multi-junction evaluation).} We report multi-seed evaluations on a $5$-junction corridor and a $16$-junction grid under SUMO's sublane model~\cite{lopez2018_sumo} with a motorcycle-dominant
  ($\approx 73\%$) fleet calibrated to Hanoi field data~\cite{ngoc2021_fivecities,cao_sano_hanoi} and time-varying demand. On the $16$-junction grid the policy improves on every camera-deployable baseline---Webster~\cite{webster1958}, max-pressure control~\cite{varaiya2013_maxpressure}, SOTL~\cite{cools2013_sotl}, and independent PPO~\cite{dewitt2020_ippo}---while remaining statistically indistinguishable from a loop-instrumented gap-actuated controller that exploits per-second in-road sensing unavailable to a camera-only deployment; on the $5$-junction corridor, classical progression control remains superior, and we analyze this reversal as a topology-dependence result that delineates where learned camera-only control pays (Section~\ref{sec:topology}). A sublane-versus-lane cross-evaluation furthermore shows that ignoring mixed-traffic dynamics materially alters policy outcomes.
\end{itemize}
We emphasize that the MAPPO algorithm itself is standard and is not claimed as a methodological contribution; it is the vehicle for evaluating the proposed camera-constrained state and reward formulation.

\subsection{Paper Organization}
The remainder of this paper is organized as follows. Section~II reviews related work and positions our contribution. Section~III formalizes the Dec-POMDP, the camera-computable observation projection, and the feature-estimator validation. Section~IV details the queue-centric vision-proxy reward, the MAPPO architecture, and the sensing-noise model. Section~V describes the experimental setup, networks, demand, baselines, and statistical protocol. Section~VI presents the main comparison, the information-restriction cost, robustness under sensing noise, generalization, and ablations. Section~VII discusses limitations, and Section~VIII concludes.